\documentclass[12pt]{article}
\usepackage[utf8]{inputenc}
\usepackage[T1]{fontenc}
\usepackage{hyperref}
\title{Yahuna Language Resources for K-12 Curriculum Development}
\author{Compiled by Claude Sonnet 5.5 (claude-sonnet-5-5) for Andr\'as Kornai}
\date{}
\begin{document}
\maketitle

\section{Introduction}
Yahuna (Yauna, Ja\'una; ISO 639-3: ynu; Glottocode yahu1241) is a Tucanoan language of Colombia. Wikipedia lists it as extinct, with an ethnic population of about 110 (2018, citing Ethnologue). Its primary documentation appears to be 19th- and early 20th-century German sources, so little is documented. Yahuna should not be confused with Yucuna, an Arawakan language with a living community, or with the Macuna people listed on a Joshua Project page titled ``Yauna''.

\section{Primary Educational Resources}
None found. No primers, textbooks or school programmes for Yahuna appeared in the searches.

\section{Linguistic and Reference Materials}
\begin{itemize}
\item Wikipedia, ``Yahuna language'': summary of status and sources. \url{https://en.wikipedia.org/wiki/Yahuna_language}
\item Wikipedia, ``Tucanoan languages'' (a mirrored copy): family context for comparison with related living languages. \url{https://static.hlt.bme.hu/semantics/external/pages/jokuc/en.wikipedia.org/wiki/Tucanoan_languages.html}
\end{itemize}
No grammar or dictionary of Yahuna was found. Any revitalization would have to use the old wordlists cited by Wikipedia (not located here) and related Tucanoan languages.

\section{Religious and Community Resources}
None found for Yahuna. The Joshua Project page \url{https://joshuaproject.net/people_groups/15948/CO} is titled ``Yauna in Colombia'', but it lists Macuna as the primary language and so is not evidence about ynu.

\section{Digital and Audio Resources}
None found.

\section{Curriculum Development Context}
There are no known speakers and no usable documentation found in the searches. A curriculum would be heritage-awareness content, such as the history of the Yahuna people and their place among Tucanoan peoples. Language teaching would need archival research first, for example into the historical wordlists. Descendant communities would have to be consulted, and no contact information was found.

\end{document}